% from ega0-19-fr.tex
\oldpage[0\textsubscript{IV-1}]{76}

à gauche, il en est de même de
$u_{00}:M/\mathfrak J_0M\to N/\mathfrak J_0M$ obtenu par tensorisation de
$u_0$ avec $1_{A/\mathfrak J_0}$, puisque l'on a
\[
  M/\mathfrak J_0M=(M/\mathfrak JM)\otimes_{A/\mathfrak J}(A/\mathfrak J_0).
\]
Supposons donc $\mathfrak J$ maximal, de sorte que $A/\mathfrak J$ est un
corps. Il suffit évidemment de montrer que $M$ est un $A$-module \emph{libre}
sous les conditions de l'énoncé~: en effet, $\det(w_0)$ est alors l'image
canonique de $\det(w)$ dans $A/\mathfrak J$, donc $\det(w)$ n'appartient pas à
l'idéal $\mathfrak J$ et est par suite inversible. Or le
$(A/\mathfrak J)$-espace vectoriel $M/\mathfrak JM$ étant libre de type fini,
il y a un $A$-module de type fini $L$ et un $A$-homomorphisme $f:L\to M$ tel
que l'homomorphisme $f_0:L/\mathfrak JL\to M/\mathfrak JM$ déduit de $f$ par
passage aux quotients soit bijectif. Comme $M$ est de type fini, on en conclut
tout d'abord que $f$ est surjectif par le lemme de Nakayama (Bourbaki,
\emph{Alg. comm.}, chap.~II, §~3, n\textsuperscript{o}~2, cor.~1 de la
prop.~4)~; en outre, si $g=u\circ f$, l'homomorphisme
$g_0:L/\mathfrak JL\to N/\mathfrak JN$ déduit par passage aux quotients est
inversible à gauche, et puisqu'ici $L$ est libre, la remarque du début prouve
que $g$ est lui-même inversible à gauche~; or cela entraîne évidemment que
$f$ est \emph{injectif}, ce qui achève la démonstration.
\end{proof}

Mentionnons en passant les corollaires suivants de (19.1.10)~:

\begin{corollary}[19.1.11]
\label{0.19.1.11-fr}
Soient $A$ un anneau local, $k$ son corps résiduel, $M$ un $A$-module de type
fini, $N$ un $A$-module projectif, $u:M\to N$ un homomorphisme. Pour que $u$
soit inversible à gauche, il faut et il suffit qu'il existe un système de
générateurs $(x_i)_{1\leq i\leq m}$ de $M$ tel que les images par
$u\otimes 1:M\otimes_Ak\to N\otimes_Ak$ des $x_i\otimes1$ soient linéairement
indépendantes dans $N\otimes_Ak$~; les $x_i$ forment alors une base de $M$.
\end{corollary}

La condition est évidemment nécessaire, car si $u$ est inversible à gauche,
$M$ est un $A$-module projectif de type fini, donc libre. Inversement, si la
condition est satisfaite, il est clair que les $x_i\otimes1$ forment une base
de $M\otimes_Ak$ et que $u\otimes1$ est inversible à gauche~; il suffit donc
d'appliquer (19.1.10) à l'idéal maximal $\mathfrak J$ de $A$.

\begin{corollary}[19.1.12]
\label{0.19.1.12-fr}
Soient $A$ un anneau, $M$ un $A$-module de type fini, $N$ un $A$-module
projectif, $u:M\to N$ un homomorphisme. Pour tout idéal premier
$\mathfrak p\in\operatorname{Spec}(A)$, les conditions suivantes sont
équivalentes~:
\begin{enumerate}
  \item[a)] L'homomorphisme
  $u_{\mathfrak p}:M_{\mathfrak p}\to N_{\mathfrak p}$ est inversible à
  gauche.
  \item[b)] L'homomorphisme
  $u\otimes1:M\otimes_Ak(\mathfrak p)\to N\otimes_Ak(\mathfrak p)$ est
  injectif (on rappelle que $k(\mathfrak p)$ désigne le corps résiduel de $A$
  en $\mathfrak p$).
  \item[c)] Il existe un système fini d'éléments $x_i\in M$
  $(1\leq i\leq m)$ tel que les images des $x_i$ dans $M_{\mathfrak p}$
  engendrent $M_{\mathfrak p}$, et un système de $m$ formes linéaires
  $y_i^*$ sur $N$ $(1\leq i\leq m)$ tel que
  $\det(\langle y_i^*,u(x_j)\rangle)\notin\mathfrak p$.
  \item[d)] Il existe $f\in A-\mathfrak p$ tel que l'homomorphisme
  $u_f:M_f\to N_f$ soit inversible à gauche.
\end{enumerate}
En outre, l'ensemble des $\mathfrak p\in\operatorname{Spec}(A)$ vérifiant ces
conditions est ouvert dans $\operatorname{Spec}(A)$.
\end{corollary}

La dernière condition est conséquence triviale de d). Comme $N$ est
projectif, il est facteur direct d'un $A$-module libre $A^{(I)}$~; en outre,
comme $M$ est de type fini, $u(M)$ est contenu dans un sous-module de $A^{(I)}$
de la forme $A^n$ pour $n$ fini~; comme chacun des énoncés a), b), c), d) est
équivalent à l'énoncé correspondant où l'on remplace $N$ par $N\oplus P$ ($u$
étant toujours supposé appliquer $M$ dans $N$), on est ramené au cas où $N$
est \emph{libre de type fini}. Il est trivial que d) entraîne a), et a) et b)
sont équivalentes en vertu de (19.1.11)~; d'ailleurs a) entraîne que
$M_{\mathfrak p}$ est libre (19.1.11), donc a) entraîne c), en prenant pour les
$x_i$ une famille dont les images dans $M_{\mathfrak p}$ forment une base de ce

